\section{Actors and eventual references}
\label{sec:actors}

This section lifts the sequential relation of Section~\ref{sec:machine} to a
system of isolated, single-threaded actors.  The lift does not make an
asynchronous send an ordinary send in disguise: evaluation returns a promise
immediately, the request becomes a mailbox item, and ordinary dispatch begins
only when the target actor starts a later turn.

\subsection{Global domains and isolation}

Let the following identity sets be pairwise disjoint and disjoint from those in
Equation~\eqref{eq:runtime-identities}:
\begin{equation}\label{eq:eventual-identities}
 p\in\mathit{PromiseId},\qquad r\in\mathit{FarRefId},\qquad
 u\in\mathit{EventId},\qquad
 \mathit{Ref}=\mathit{ObjRef}\uplus\mathit{PromiseId}\uplus
                 \mathit{FarRefId}.
\end{equation}
Promises and far references are Newspeak objects with platform-supplied
classes, but their routing state is recorded globally below.  An ordinary send
to such a handle invokes its local platform protocol; only
$\kw{asyncSend}$ dereferences it for remote delivery.

An actor has a private heap $H_A$.  A platform may additionally maintain a
value heap $V$ shared by actors managed by one virtual machine.  The heap visible to the
sequential machine running actor $A$ is
\begin{equation}\label{eq:actor-heap-view}
 \widehat H_A=V\uplus H_A,\qquad
 A\ne B\Longrightarrow
 \dom(H_A)\cap\dom(H_B)=\varnothing,\qquad
 \dom(V)\cap\dom(H_A)=\varnothing.
\end{equation}
Only objects satisfying the value predicate of Section~\ref{sec:modules} may
enter $V$.  A distributed implementation may instead copy their complete
immutable reachable graph.  We write
\begin{equation}\label{eq:value-transfer-contract}
 \Omega\vdash\mathsf{valTransfer}_{A\Rightarrow B}(x)
   \Downarrow\langle\Omega',y\rangle
\end{equation}
for either permitted choice; it is defined only for a fully initialized value
$x$, preserves its class and immutable reachable graph up to identity
isomorphism, and neither aliases nor modifies mutable actor-local state.
Sharing is the identity case.  The reflective layer will distinguish actors in
one VM, which may share values and observe common code change, from actors in
different VMs, which must use the copying case.

A far-reference table and promise table have the forms
\begin{equation}\label{eq:eventual-tables}
\begin{aligned}
 \Phi(r)&=\langle A,o\rangle,\\
 \Pi(p)&::=\kw{pending}(A,W)
       \mid\kw{fulfilled}(A,x)\mid\kw{broken}(A,x),\\
 W&::=\epsilon\mid W\mathbin{\cdot}
       \langle B,\mu,p_r,\kappa\rangle.
\end{aligned}
\end{equation}
$\Phi(r)=\langle A,o\rangle$ means that $r$ designates $o$ in $H_A$.
The actor in a promise state is the actor relative to whose heap its result is
represented.  Pending sends are retained in registration order; a waiter
$\langle B,\mu,p_r,\kappa\rangle$ requests that, after settlement, actor $B$
route $\mu$ to the result and settle $p_r$.  The causal history $\kappa$ is
defined below.

Remote representation is an actor-relative relation
$\Omega\vdash\mathsf{rr}_{A\Rightarrow B}(x)
\Downarrow\langle\Omega',y\rangle$.  It is the least relation given by the
following rules:
\[
\dfrac{A=B}
 {\Omega\vdash\mathsf{rr}_{A\Rightarrow B}(x)
   \Downarrow\langle\Omega,x\rangle}
\quad(\rulelabel{RR-Same})
\]
\[
\dfrac{A\ne B\quad \mathsf{value}_{P,\widehat H_A}(x)\quad
       \Omega\vdash\mathsf{valTransfer}_{A\Rightarrow B}(x)
         \Downarrow\langle\Omega',y\rangle}
 {\Omega\vdash\mathsf{rr}_{A\Rightarrow B}(x)
   \Downarrow\langle\Omega',y\rangle}
\quad(\rulelabel{RR-Value})
\]
\[
\dfrac{A\ne B\quad \Phi(r)=\langle B,o\rangle}
 {\Omega\vdash\mathsf{rr}_{A\Rightarrow B}(r)
   \Downarrow\langle\Omega,o\rangle}
\quad(\rulelabel{RR-Far-Home})
\qquad
\dfrac{A\ne B\quad \Phi(r)=\langle C,o\rangle\quad C\ne B}
 {\Omega\vdash\mathsf{rr}_{A\Rightarrow B}(r)
   \Downarrow\langle\Omega,r\rangle}
\quad(\rulelabel{RR-Far-Pass})
\]
\[
\dfrac{A\ne B\quad p\in\dom(\Pi)}
 {\Omega\vdash\mathsf{rr}_{A\Rightarrow B}(p)
   \Downarrow\langle\Omega,p\rangle}
\quad(\rulelabel{RR-Promise})
\]
\[
\dfrac{A\ne B\quad o\in\dom(H_A)\quad
       \neg\mathsf{value}_{P,\widehat H_A}(o)\quad r\text{ fresh}\quad
       \Omega'=\Omega[\Phi(r)\mapsto\langle A,o\rangle]}
 {\Omega\vdash\mathsf{rr}_{A\Rightarrow B}(o)
   \Downarrow\langle\Omega',r\rangle}
\quad(\rulelabel{RR-Near-Far})
\]
Thus a far reference passed back to its target actor becomes the near referent,
and a far reference passed to a third actor does not acquire a proxy chain.
Promise identities are monotonic eventual capabilities and likewise pass
without a proxy layer.  Applying remote representation componentwise defines
message transport:
\begin{equation}\label{eq:message-transfer}
 \dfrac{\Omega_0=\Omega\quad
   \Omega_{i-1}\vdash\mathsf{rr}_{A\Rightarrow B}(o_i)
      \Downarrow\langle\Omega_i,o_i'\rangle\ (1\le i\le n)}
 {\Omega\vdash\mathsf{msgRR}_{A\Rightarrow B}
      (\langle s,\langle o_1,\ldots,o_n\rangle\rangle)
   \Downarrow
      \langle\Omega_n,\langle s,\langle o_1',\ldots,o_n'\rangle\rangle\rangle}.
\end{equation}
The selector is a symbol value and is therefore unchanged.

\subsection{Actor configurations and causal delivery}

A mailbox or network packet is
\begin{equation}\label{eq:packet-grammar}
\begin{split}
 b::={}&\kw{app}\langle u,A,B,o,\mu,p,\kappa\rangle\\
  \mid{}&\kw{settlementPacket}\langle u,A,B,p,\chi,x,\kappa\rangle\\
  \mid{}&\kw{wake}\langle u,A,B,p,p_r,\chi,x,\mu,\kappa_w,\kappa\rangle,
 \qquad \chi\in\{\kw{ok},\kw{error}\}.
\end{split}
\end{equation}
The first two actor fields are source and destination.  Every packet has a
fresh event identity $u$ and carries the source actor's causally prior event
set $\kappa$.  A wake names both the settled promise $p$ that caused the wake
and the dependent result promise $p_r$; it also retains $\kappa_w$, the
history at which its dependent send was originally issued.

An actor is idle, is executing one sequential configuration on behalf of an
application packet, or is suspended at a debugger safepoint:
\begin{equation}\label{eq:actor-run-state}
 \begin{aligned}
 \mathcal X(A)::={}&\kw{idle}\mid
   \kw{runningTurn}\langle\mathcal A,t,p,u\rangle\mid\\
 &\kw{pausedTurn}\langle\mathcal A,t,p,u,r_s,d_s\rangle.
 \end{aligned}
\end{equation}
Here $d_s\in\mathit{StopReason}$.  The ordinary actor rules have no premise
for \kw{pausedTurn}; deliveries may still append to its mailbox, but the actor
executes no turn step until an authorized debugger operation resumes it.
Section~\ref{sec:debugging} defines suspension, stack replacement, and
resumption.
The global state is
\begin{equation}\label{eq:actor-configuration}
 \Omega=\langle P,V,\mathcal H,\mathcal X,\mathcal Q,\mathcal N,
              \Pi,\Phi,\Gamma,\mathcal E,\mathcal D,w,n\rangle,
 \qquad w\in\mathit{VMId},\quad n\in\mathbb N.
\end{equation}
$w$ identifies the VM cohort and $n$ is its committed program version.  Every
actor in $\dom(\mathcal H)$ is managed by $w$ and observes the same $P$ and
$V$.  Section~\ref{sec:reflection} changes $P$ and $n$ atomically.  A system of
multiple VMs is a map from distinct $w$ to configurations of this form; its
program components evolve independently.
$\Gamma_\Omega$, $\Pi_\Omega$, and analogous subscripts name a component of a
particular configuration when more than one configuration occurs in a rule;
the subscript is omitted otherwise.
$\mathcal Q(A)$ is a FIFO mailbox, $\mathcal N$ is the multiset of undelivered
cross-actor packets, and $\Gamma(A)$ is a finite downward-closed causal
history.  The permanent event ledger records
\begin{equation}\label{eq:event-ledger}
 \mathcal E(u)=\langle A,B,\kappa\rangle,\qquad
 \mathsf{dst}_{\Omega}(u)=B,\qquad
 \mathcal D(B)\in\mathit{EventId}^{*}.
\end{equation}
$\mathcal D(B)$ lists, in arrival order, events delivered to $B$; a failed
network event is marked retired at the same position.  Processed event records
are retained because later causal histories may still mention them.

All application and system packets are emitted by the helper
\begin{equation}\label{eq:emit-packet}
\begin{split}
 &\mathsf{emit}(\Omega,A,B,\beta,\kappa)=\Omega_2,\\
 &u\text{ fresh},\quad b=\beta(u,A,B,\kappa),\quad
 \Omega_1=\Omega[\mathcal E(u)\mapsto\langle A,B,\kappa\rangle,
                  \Gamma(A)\mapsto\Gamma(A)\cup\{u\}],\\
 &\Omega_2=
 \begin{cases}
  \Omega_1[\mathcal Q(B)\mapsto\mathcal Q(B)\cdot b,
           \mathcal D(B)\mapsto\mathcal D(B)\cdot u] & A=B,\\
  \Omega_1[\mathcal N\mapsto\mathcal N\uplus\{b\}] & A\ne B.
 \end{cases}
\end{split}
\end{equation}
Here $\beta$ is a packet template supplying every field except
$u,A,B,\kappa$.  A same-actor asynchronous send arrives immediately at the
tail of its own mailbox, but cannot start until the current turn finishes.

A network packet $b$ with identity $u$, destination $B$, and history $\kappa$
is eligible precisely when
\begin{equation}\label{eq:e-order-eligibility}
 \mathsf{eligible}_{\Omega}(b)\quad\Longleftrightarrow\quad
 \{u'\in\kappa\mid\mathsf{dst}_{\Omega}(u')=B\}
       \subseteq\mathsf{set}(\mathcal D(B)).
\end{equation}
An eligible packet may arrive:
\[
\dfrac{b\in\mathcal N\quad \mathsf{id}(b)=u\quad
       \mathsf{destination}(b)=B\quad
       \mathsf{eligible}_{\Omega}(b)}
 {\Omega\longrightarrow_{\mathsf{act}}
  \Omega[\mathcal N\mapsto\mathcal N\setminus\{b\},
         \mathcal Q(B)\mapsto\mathcal Q(B)\cdot b,
         \mathcal D(B)\mapsto\mathcal D(B)\cdot u]}
\quad(\rulelabel{Network-Deliver})
\]
Because mailboxes are FIFO and an actor runs at most one turn, delivery
in this order implies processing in this order.

Equation~\eqref{eq:e-order-eligibility} is a constructive form of E-order.
After $A_1$ emits $u_1$ to $A_2$, $u_1\in\Gamma(A_1)$, so a later $u_2$ from
$A_1$ to $A_2$ carries $u_1$ and cannot arrive first.  If the later packet
instead passes a far reference for $A_2$ to $A_3$, delivery merges its history
into $\Gamma(A_3)$; a subsequent send by $A_3$ to that referent therefore also
carries $u_1$ and cannot arrive at $A_2$ first.  These are exactly the two
ordering clauses in Specification \S5.6, generalized by transitive causal
history.

Maximal executions are required to be weakly fair:
\begin{equation}\label{eq:actor-fairness}
 \mathsf{continuouslyEnabled}(\alpha,i)
   \Longrightarrow\exists j\ge i.\ \mathsf{taken}(\alpha,j),
\end{equation}
for packet delivery, dequeue of the head of a nonempty idle mailbox, and an
enabled step of a running actor.  Fairness does not rescue an actor whose
current turn diverges; its later mailbox entries may consequently remain
unprocessed.

\subsection{Asynchronous-send evaluation and routing}

Extend the request and frame grammars in
Equations~\eqref{eq:request-grammar} and~\eqref{eq:frame-grammar} by
\begin{equation}\label{eq:async-frames}
 \theta::=\cdots\mid\kw{eventualRequest}(x),\qquad
 F::=\cdots\mid\kw{asyncRecv}(s,\bar e).
\end{equation}
The receiver is evaluated first:
\[
 \langle P,H,\mathcal A\cdot\langle a,\pi\rangle,
        \kw{asyncSend}(e,s,\bar e)\rangle
 \longmapsto
 \langle P,H,\mathcal A\cdot
        \langle a,\pi\cdot\kw{asyncRecv}(s,\bar e)\rangle,e\rangle
 \quad(\rulelabel{Async-Receiver}).
\]
After it yields $x$, the existing argument frames of
Equation~\eqref{eq:argument-continuations} evaluate the arguments from left to
right:
\begin{equation}\label{eq:async-receiver-continuations}
\begin{array}{c}
 \langle\pi\cdot\kw{asyncRecv}(s,\epsilon),x\rangle
  \longmapsto
 \langle\pi,\kw{dispatch}(\kw{eventualRequest}(x),
                          \langle s,\epsilon\rangle)\rangle,\\[2mm]
 \langle\pi\cdot\kw{asyncRecv}(s,\langle e_1,\bar e\rangle),x\rangle
  \longmapsto
 \langle\pi\cdot\kw{args}(\kw{eventualRequest}(x),s,\epsilon,\bar e),e_1\rangle.
\end{array}
\end{equation}

For a non-promise receiver, define its target by
\begin{equation}\label{eq:eventual-target}
 \mathsf{target}_{\Omega,A}(x)=
 \begin{cases}
  \langle A,x\rangle & x\in\dom(\widehat H_A),\\
  \langle B,o\rangle & \Phi(x)=\langle B,o\rangle.
 \end{cases}
\end{equation}
The routing judgment
$\Omega,A\vdash\mathsf{route}(x,\mu,p_r,\kappa)\Downarrow\Omega'$ is
defined by the following four cases.  A known target transports the arguments
and emits an application packet:
{\small
\[
\dfrac{\begin{gathered}
 \mathsf{target}_{\Omega,A}(x)=\langle B,o\rangle\\
 \Omega\vdash\mathsf{msgRR}_{A\Rightarrow B}(\mu)
    \Downarrow\langle\Omega_1,\mu'\rangle\\
 \Omega'=\mathsf{emit}(\Omega_1,A,B,
    \lambda u,A,B,\kappa.\kw{app}\langle
       u,A,B,o,\mu',p_r,\kappa\rangle,\kappa)
\end{gathered}}
 {\Omega,A\vdash\mathsf{route}(x,\mu,p_r,\kappa)\Downarrow\Omega'}
\quad(\rulelabel{Route-Target})
\]}
An unresolved promise records the send without blocking its actor:
\[
\dfrac{\Pi(p)=\kw{pending}(C,W)\quad
       \Omega'=\Omega[\Pi(p)\mapsto
          \kw{pending}(C,W\cdot\langle A,\mu,p_r,\kappa\rangle)]}
 {\Omega,A\vdash\mathsf{route}(p,\mu,p_r,\kappa)\Downarrow\Omega'}
\quad(\rulelabel{Route-Pending})
\]
For a fulfilled promise, its result is first represented in the requesting
actor and routing repeats:
{\small\[
\dfrac{\Pi(p)=\kw{fulfilled}(C,x)\quad
       \Omega\vdash\mathsf{rr}_{C\Rightarrow A}(x)
          \Downarrow\langle\Omega_1,y\rangle\quad
       \Omega_1,A\vdash\mathsf{route}(y,\mu,p_r,\kappa)
          \Downarrow\Omega'}
 {\Omega,A\vdash\mathsf{route}(p,\mu,p_r,\kappa)\Downarrow\Omega'}
\quad(\rulelabel{Route-Fulfilled})
\]}
For a broken promise, the result promise is broken with the same represented
exception:
{\small\[
\dfrac{\begin{gathered}\Pi(p)=\kw{broken}(C,x)\\
       \Omega\vdash\mathsf{rr}_{C\Rightarrow A}(x)
          \Downarrow\langle\Omega_1,y\rangle\\
       \Omega_1,A\vdash\mathsf{settlePromise}(p_r,\kw{error},A,y,\kappa)
          \Downarrow\Omega'\end{gathered}}
 {\Omega,A\vdash\mathsf{route}(p,\mu,p_r,\kappa)\Downarrow\Omega'}
\quad(\rulelabel{Route-Broken})
\]}
The last rule makes explicit the otherwise unspecified meaning of sends to a
broken promise: break propagation is the conventional eventual-send behavior,
but is recorded as a language-design choice in the coverage ledger.

Routing is an actor-level action.  The asynchronous expression itself
allocates its result promise and immediately leaves that promise as the value
of the expression:
{\small
\[
\dfrac{\begin{gathered}
 \mathcal X(A)=\kw{runningTurn}\langle\mathcal A\cdot\langle a,\pi\rangle,
   \kw{dispatch}(\kw{eventualRequest}(x),\mu),p_0,u_0\rangle\\
 p_r\text{ fresh}\qquad
 \Omega_0=\Omega[\Pi(p_r)\mapsto\kw{pending}(A,\epsilon)]\\
 \Omega_0,A\vdash\mathsf{route}(x,\mu,p_r,\Gamma(A))\Downarrow\Omega_1
\end{gathered}}
 {\begin{aligned}
  \Omega\longrightarrow_{\mathsf{act}}\Omega_1[\mathcal X(A)\mapsto{}&
   \kw{runningTurn}\langle\mathcal A\cdot\langle a,\pi\rangle,\\[-1mm]
   &p_r,p_0,u_0\rangle]
 \end{aligned}}
\quad(\rulelabel{Async-Return})
\]}
The displayed run record in the conclusion abbreviates replacement of its
control component by $p_r$; its reply promise and triggering event remain
$p_0,u_0$.

\subsection{Turns and promise settlement}

When an application packet reaches the head of an idle actor's mailbox, it
starts a turn.  With fresh top activation $a_T$ and
$H' =\mathsf{newTop}(\widehat H_B,a_T)$ from
Equation~\eqref{eq:new-top-activation}:
{\small
\[
\dfrac{\mathcal X(B)=\kw{idle}\quad
 \mathcal Q(B)=\kw{app}\langle u,A,B,o,\mu,p,\kappa\rangle\cdot Q}
 {\begin{aligned}
  \Omega\longrightarrow_{\mathsf{act}}\Omega[&\mathcal X(B)\mapsto
    \kw{runningTurn}\langle\epsilon\cdot\langle a_T,\epsilon\rangle,
       \kw{dispatch}(\kw{ordinaryRequest}(o),\mu),p,u\rangle,\\[-1mm]
   &\mathcal Q(B)\mapsto Q,
    \Gamma(B)\mapsto\Gamma(B)\cup\kappa\cup\{u\},
    \mathcal H(B)\mapsto H'\setminus V]
 \end{aligned}}
\quad(\rulelabel{Turn-Start})
\]}
While $B$ is running, each applicable sequential step is lifted pointwise:
\[
\dfrac{\langle P,\widehat H_B,\mathcal A,t\rangle
          \longmapsto\langle P,H',\mathcal A',t'\rangle}
 {\Omega\longrightarrow_{\mathsf{act}}
  \Omega[\mathcal H(B)\mapsto H'\setminus V,
         \mathcal X(B)\mapsto
          \kw{runningTurn}\langle\mathcal A',t',p,u\rangle]}
\quad(\rulelabel{Actor-Sequential})
\]
It is inapplicable to the eventual dispatch intercepted by
\rulelabel{Async-Return}.  No second packet can start in $B$ until this turn
ends.

Settlement wakes every recorded dependent send.  Define the ordered helper
{\small
\begin{equation}\label{eq:wake-all}
\begin{aligned}
 \mathsf{wakeAll}(\Omega,B,p,\chi,y,\epsilon)&=\Omega,\\
 \mathsf{wakeAll}(\Omega,B,p,\chi,y,
    \langle C,\mu,p_r,\kappa_w\rangle\cdot W)
   &=\mathsf{wakeAll}(\Omega_1,B,p,\chi,y,W),\\
 \Omega_1&=\mathsf{emit}(\Omega,B,C,
   \lambda u,B,C,\kappa.\kw{wake}\langle
      u,B,C,p,p_r,\chi,y,\mu,\kappa_w,\kappa\rangle,\Gamma(B)).
\end{aligned}
\end{equation}}
Earlier waiters for the same actor are emitted first and therefore ordered by
Equation~\eqref{eq:e-order-eligibility}.  Promise settlement is the partial
judgment
{\small
\begin{equation}\label{eq:promise-settlement}
\dfrac{\begin{gathered}
 \Pi(p)=\kw{pending}(B,W)\qquad \kappa\subseteq\Gamma(B)\\
 \Omega\vdash\mathsf{rr}_{A\Rightarrow B}(x)
    \Downarrow\langle\Omega_0,y\rangle\\
 \Omega_1=
 \begin{cases}
  \Omega_0[\Pi(p)\mapsto\kw{fulfilled}(B,y)]&\chi=\kw{ok},\\
  \Omega_0[\Pi(p)\mapsto\kw{broken}(B,y)]&\chi=\kw{error}
 \end{cases}
 \qquad \Omega'=\mathsf{wakeAll}(\Omega_1,B,p,\chi,y,W)
\end{gathered}}
 {\Omega,B\vdash\mathsf{settlePromise}(p,\chi,A,x,\kappa)
    \Downarrow\Omega'}.
\end{equation}}
There is no rule for settling an already settled promise, so fulfillment and
breaking are mutually exclusive and occur at most once.

When a turn returns normally, a promise owned by the same actor settles
immediately; a remote owner's settlement is sent subsequently:
{\small\[
\dfrac{\begin{gathered}
 \mathcal X(B)=\kw{runningTurn}\langle\epsilon,\kw{halt}(v),p,u\rangle
 \qquad \Pi(p)=\kw{pending}(B,W)\\
 \Omega,B\vdash\mathsf{settlePromise}(p,\kw{ok},B,v,\Gamma(B))
    \Downarrow\Omega'
\end{gathered}}
 {\Omega\longrightarrow_{\mathsf{act}}
  \Omega'[\mathcal X(B)\mapsto\kw{idle}]}
\quad(\rulelabel{Turn-Fulfill-Local})
\]}
{\small\[
\dfrac{\begin{gathered}
 \mathcal X(B)=\kw{runningTurn}\langle\epsilon,\kw{halt}(v),p,u\rangle\\
 \Pi(p)=\kw{pending}(A,W)\qquad A\ne B\\
 \Omega'=\mathsf{emit}(\Omega,B,A,
  \lambda u,B,A,\kappa.\kw{settlementPacket}\langle
     u,B,A,p,\kw{ok},v,\kappa\rangle,\Gamma(B))
\end{gathered}}
 {\Omega\longrightarrow_{\mathsf{act}}
  \Omega'[\mathcal X(B)\mapsto\kw{idle}]}
\quad(\rulelabel{Turn-Fulfill-Remote})
\]}
If exception signaling escapes the turn's top activation, the
exception-library boundary produces $\kw{abort}(x)$ instead of
$\kw{halt}(v)$.  A local owner's promise is broken immediately:
{\small\[
\dfrac{\begin{gathered}
 \mathcal X(B)=\kw{runningTurn}\langle\epsilon,\kw{abort}(x),p,u\rangle
 \qquad \Pi(p)=\kw{pending}(B,W)\\
 \Omega,B\vdash\mathsf{settlePromise}(p,\kw{error},B,x,\Gamma(B))
    \Downarrow\Omega'
\end{gathered}}
 {\Omega\longrightarrow_{\mathsf{act}}
  \Omega'[\mathcal X(B)\mapsto\kw{idle}]}
\quad(\rulelabel{Turn-Break-Local})
\]}
A remote owner receives an error settlement subsequently:
{\small\[
\dfrac{\begin{gathered}
 \mathcal X(B)=\kw{runningTurn}\langle\epsilon,\kw{abort}(x),p,u\rangle\\
 \Pi(p)=\kw{pending}(A,W)\qquad A\ne B\\
 \Omega'=\mathsf{emit}(\Omega,B,A,
  \lambda u,B,A,\kappa.\kw{settlementPacket}\langle
     u,B,A,p,\kw{error},x,\kappa\rangle,\Gamma(B))
\end{gathered}}
 {\Omega\longrightarrow_{\mathsf{act}}
  \Omega'[\mathcal X(B)\mapsto\kw{idle}]}
\quad(\rulelabel{Turn-Break-Remote})
\]}
These rules are the only additional obligation placed on the resumable
exception library by actors.

An idle promise owner processes a settlement packet at the head of its FIFO
mailbox by merging the packet history and applying
Equation~\eqref{eq:promise-settlement}:
{\small\[
\dfrac{\begin{gathered}
 \mathcal X(B)=\kw{idle}\qquad
 \mathcal Q(B)=\kw{settlementPacket}\langle
    u,A,B,p,\chi,x,\kappa\rangle\cdot Q\\
 \Omega_0=\Omega[\mathcal Q(B)\mapsto Q,
                  \Gamma(B)\mapsto\Gamma(B)\cup\kappa\cup\{u\}]\\
 \Omega_0,B\vdash\mathsf{settlePromise}(p,\chi,A,x,\Gamma_0(B))
    \Downarrow\Omega'
\end{gathered}}
 {\Omega\longrightarrow_{\mathsf{act}}\Omega'}
\quad(\rulelabel{Settlement-Packet-Dequeue})
\]}
A fulfilled wake represents the result in the waiting actor and resumes the
stored routing operation; a broken wake settles the dependent promise as
broken:
{\small\[
\dfrac{\begin{gathered}
 \mathcal X(C)=\kw{idle}\qquad
 \mathcal Q(C)=\kw{wake}\langle
  u,B,C,p,p_r,\kw{ok},x,\mu,\kappa_w,\kappa\rangle\cdot Q\\
 \Omega_0=\Omega[\mathcal Q(C)\mapsto Q,
   \Gamma(C)\mapsto\Gamma(C)\cup\kappa\cup\kappa_w\cup\{u\}]\\
 \Omega_0\vdash\mathsf{rr}_{B\Rightarrow C}(x)
    \Downarrow\langle\Omega_1,y\rangle\\
 \Omega_1,C\vdash\mathsf{route}(y,\mu,p_r,\Gamma_1(C))
    \Downarrow\Omega'
\end{gathered}}
 {\Omega\longrightarrow_{\mathsf{act}}\Omega'}
\quad(\rulelabel{Wake-Fulfilled})
\]}
{\small\[
\dfrac{\begin{gathered}
 \mathcal X(C)=\kw{idle}\qquad
 \mathcal Q(C)=\kw{wake}\langle
  u,B,C,p,p_r,\kw{error},x,\mu,\kappa_w,\kappa\rangle\cdot Q\\
 \Omega_0=\Omega[\mathcal Q(C)\mapsto Q,
   \Gamma(C)\mapsto\Gamma(C)\cup\kappa\cup\kappa_w\cup\{u\}]\\
 \Omega_0,C\vdash\mathsf{settlePromise}(p_r,\kw{error},B,x,\Gamma_0(C))
    \Downarrow\Omega'
\end{gathered}}
 {\Omega\longrightarrow_{\mathsf{act}}\Omega'}
\quad(\rulelabel{Wake-Broken})
\]}
These system packets are handled only between user turns, so remote settlement
cannot mutate an actor's heap or promise observations during an executing turn.
A communication subsystem may retire an undelivered application event and
emit an error settlement instead of delivering it.  Its failure object must
satisfy the same remote-representation boundary, the event is entered in
$\mathcal D$ as retired so it cannot block causal successors, and delivery and
failure of the same event are mutually exclusive.

\subsection{Actor creation}

Actor creation is a platform primitive rather than a new surface expression.
Elaboration supplies, for a top-level mixin, the class-side mixin and factory
descriptor needed to apply it to \Object:
\begin{equation}\label{eq:actor-seed}
 P.\mathit{actorSeed}(M)=\langle D,M_c,q\rangle
 \Longrightarrow
 \decl_P(M)=D\ \land\ P.\mathit{owner}(D)=\kw{topOwner}.
\end{equation}
Let $\mathsf{base}_P(B)$ be a fresh actor's platform heap, containing the
distinguished classes and canonical values but no user-created mutable object.
For fresh $B,C,C_m,r$, actor creation is
{\small
\begin{equation}\label{eq:create-actor}
\dfrac{\begin{gathered}
 P.\mathit{actorSeed}(M)=\langle D,M_c,q\rangle\qquad
 H_0=\mathsf{base}_P(B)\\
 H_1=\mathsf{allocClass}(H_0,C,C_m,\kw{actor}(D),M,
          P.\mathit{platform}(\kw{Object}),\nil,M_c,q)\\
 \Omega'=\Omega[\mathcal H(B)\mapsto H_1,
   \mathcal X(B)\mapsto\kw{idle},\mathcal Q(B)\mapsto\epsilon,
   \Gamma(B)\mapsto\varnothing,\\[-1mm]
 \hspace{27mm}\mathcal D(B)\mapsto\epsilon,
   \Phi(r)\mapsto\langle B,C\rangle]
\end{gathered}}
 {\Omega,A\vdash\mathsf{spawn}(M)\Downarrow\langle\Omega',r\rangle}.
\end{equation}}
The result is explicitly a far reference to the new class, as required by the
Specification.  The $\kw{actor}(D)$ provenance distinguishes that class from
evaluation of the top-level class expression itself, so it is not accidentally
classified as a module definition by Equation~\eqref{eq:module-definition}.
Its primary factory can then be invoked asynchronously in the ordinary way.
